\clearpage
\begin{explanationbox}
\textbf{第一步，只看熟悉的线段位置。}用普通点 $P$ 暂代焦点，先分清两端在它的两侧还是同侧。\begin{center}\begin{tikzpicture}[font=\small,>=stealth]
\draw[rounded corners=3pt,gray!45] (0,0) rectangle (14.2,12.2);

\node at (7.1,11.55) {先拿掉双曲线：只比较一条直线上的三个点};
\node at (7.1,10.25) {上图：中间点 $P$，两端在相反方向};
\draw[gray] (1.2,8.6)--(12.9,8.6);
\draw[blue!75!black,very thick] (6.2,8.6)--(10.2,8.6);
\draw[orange!85!black,very thick] (6.2,8.6)--(3.2,8.6);
\fill (10.2,8.6) circle(2pt) node[above=9pt] {$A$};
\fill (6.2,8.6) circle(2pt) node[above=9pt] {$P$};
\fill (3.2,8.6) circle(2pt) node[above=9pt] {$B$};
\draw[->,blue!75!black] (6.5,7.9)--(9.7,7.9);
\draw[->,orange!85!black] (5.9,7.9)--(3.7,7.9);
\node at (7.1,6.85) {$AB=PA+PB$，两段的水平位移方向相反。};
\node at (7.1,5.35) {下图：外部点 $P$，两端在同一方向};
\draw[gray] (1.2,3.8)--(12.9,3.8);
\draw[orange!85!black,very thick] (11.2,3.8)--(2.2,3.8);
\draw[blue!75!black,very thick] (11.2,3.8)--(7.2,3.8);
\fill (11.2,3.8) circle(2pt) node[above=9pt] {$P$};
\fill (7.2,3.8) circle(2pt) node[above=9pt] {$A$};
\fill (2.2,3.8) circle(2pt) node[above=9pt] {$B$};
\draw[->,blue!75!black] (10.8,3.1)--(7.7,3.1);
\draw[->,orange!85!black] (10.8,2.6)--(2.7,2.6);
\node at (7.1,1.6) {$AB=PB-PA$，两段的水平位移方向相同。};
\node at (7.1,.55) {此图只认位置和方向；倒数的加减还需要后面的焦半径关系。};
\end{tikzpicture}
\end{center}
不要把“直线经过一个点”理解成“线段一定包含这个点”。上图是两条反向射线，下图是同一条射线。
\end{explanationbox}
\clearpage
\begin{explanationbox}
\textbf{第二步，看相似三角形的投影比。}仍然没有双曲线、焦点和新公式；横向位移可以为负，点到点的长度始终为正。\begin{center}\begin{tikzpicture}[font=\small,>=stealth]
\draw[rounded corners=3pt,gray!45] (0,0) rectangle (14.2,16.2);
\node at (7.1,15.65) {再看初中相似：同向的投影比相等，反向的投影比相反};\node at (7.1,14.1) {反向：普通点 $P$ 在两端中间};
\begin{scope}[shift={(5.2,11.0)},scale=1.1]
\draw[gray] (-3,0)--(2,0);
\draw[orange!85!black,very thick] (0,0)--(-0.6000000000,-0.8000000000);\draw[blue!75!black,very thick] (0,0)--(1.2000000000,1.6000000000);
\draw[gray,dashed] (1.2000000000,1.6000000000)--(1.2,0);\fill (1.2000000000,1.6000000000) circle(2pt) node[above=8pt] {$A$};
\draw[gray,dashed] (-0.6000000000,-0.8000000000)--(-0.6,0);\fill (-0.6000000000,-0.8000000000) circle(2pt) node[below=8pt] {$B$};
\draw[gray] (1.05,0)--(1.05,.15)--(1.2,.15);\draw[gray] (-.45,0)--(-.45,-.15)--(-.6,-.15);
\fill (0,0) circle(2pt) node[below right=8pt] {$P$};\end{scope}
\node[align=left,anchor=west] at (8.6,11.9) {两个直角三角形相似\\[13pt]水平位移一正一负};
\node at (7.1,9.25) {$\displaystyle \dfrac{x_A-x_P}{PA}=-\dfrac{x_B-x_P}{PB}$};
\node at (7.1,7.85) {同向：普通点 $P$ 在两端外面};
\begin{scope}[shift={(5.2,4.75)},scale=1.1]
\draw[gray] (-3,0)--(2,0);
\draw[orange!85!black,very thick] (0,0)--(-1.8000000000,2.4000000000);\draw[blue!75!black,very thick] (0,0)--(-0.6000000000,0.8000000000);
\draw[gray,dashed] (-0.6000000000,0.8000000000)--(-0.6,0);\fill (-0.6000000000,0.8000000000) circle(2pt) node[above=8pt] {$A$};
\draw[gray,dashed] (-1.8000000000,2.4000000000)--(-1.8,0);\fill (-1.8000000000,2.4000000000) circle(2pt) node[above=8pt] {$B$};
\draw[gray] (-.45,0)--(-.45,.15)--(-.6,.15);\draw[gray] (-1.65,0)--(-1.65,.15)--(-1.8,.15);
\fill (0,0) circle(2pt) node[below right=8pt] {$P$};\end{scope}
\node[align=left,anchor=west] at (8.6,5.65) {两个直角三角形相似\\[13pt]水平位移都是负数};
\node at (7.1,3.0) {$\displaystyle \dfrac{x_A-x_P}{PA}=\dfrac{x_B-x_P}{PB}$};
\node at (7.1,1.3) {$x_P,x_A,x_B$ 分别是三点的横坐标。};\node at (7.1,.6) {所有 $PA,PB$ 都是正长度；分子才是有正负号的水平位移。};
\end{tikzpicture}
\end{center}
竖直线的水平位移都是零，比值关系仍成立；水平线虽没有非退化直角三角形，也可直接核对位移与长度的比例。
\end{explanationbox}
\clearpage
\begin{explanationbox}
\textbf{第三步，普通点 $P$ 对应右焦点 $F$。先看同支。}两点都在右支，焦点在中间；蓝段与橙段方向相反。\begin{center}\begin{tikzpicture}[font=\small,>=stealth]
\draw[rounded corners=3pt,gray!45] (0,0) rectangle (14.2,13.2);
\node at (7.1,12.55) {放回双曲线：同支、反向，两段的倒数相加};\begin{scope}[shift={(4.5,7.5)},scale=0.37]
\begin{scope}\clip (-10,-9) rectangle (10,9);
\draw[teal!75!black,thick,samples=110,domain=-1.9:1.9,variable=\t] plot ({3*cosh(\t)},{4*sinh(\t)});
\draw[teal!75!black,thick,samples=110,domain=-1.9:1.9,variable=\t] plot ({-3*cosh(\t)},{4*sinh(\t)});
\end{scope}
\draw[->,gray] (-10.2,0)--(10.3,0);
\draw[orange!85!black,very thick] (5.0,0)--(4.2000000000,-3.9191835885);\draw[blue!75!black,very thick] (5.0,0)--(6.6000000000,7.8383671769);
\fill (6.6000000000,7.8383671769) circle(6pt) node[above right=7pt] {$A$};\fill (4.2000000000,-3.9191835885) circle(6pt) node[below=8pt] {$B$};
\fill (5.0,0) circle(6pt) node[below right=9pt] {$F$};
\end{scope}

\node[anchor=west,align=left] at (9.5,8.3)
 {图取 $a=3,b=4,c=5$\\[13pt]$AF=8,\quad BF=4$\\[13pt]两点都在右支\\[13pt]$F$ 在 $AB$ 内\\[13pt]方向相反};
\node at (7.1,3.65) {$\displaystyle \frac18+\frac14=\frac38=\frac{2a}{b^2}$};
\fill[blue!30] (3,2.25) rectangle (4.4,2.75);
\fill[orange!35] (4.4,2.25) rectangle (7.2,2.75);
\draw[gray] (3,2.25) rectangle (7.2,2.75);
\node[blue!75!black] at (3.7,1.65) {$1/AF=1/8$};
\node[orange!85!black] at (5.8,1.0) {$1/BF=1/4$};
\node[anchor=west] at (8.1,2.5) {两块合起来，才是固定的 $3/8$};
\node at (7.1,.45) {条形图表示倒数；上面的蓝、橙线段表示原长度，两者不能混读。};
\end{tikzpicture}
\end{center}
令 $r_A=|AF|$、$r_B=|BF|$，$e=c/a$、$p=b^2/a$。右支焦半径关系 $r=ex-a$ 给出 $x-c=(r-p)/e$。
把它代入上一页的反向比例，便得到
\[
 \frac{r_A-p}{er_A}=-\frac{r_B-p}{er_B}
 \quad\Longrightarrow\quad \frac1{r_A}+\frac1{r_B}=\frac2p.
\]
倒数来自“水平位移除以线段长”；不是额外猜出来的形式。
\end{explanationbox}
\clearpage
\begin{explanationbox}
\textbf{第四步，再看异支。}从右焦点向同一方向走，先到右支 $A$，再到左支 $B$。焦点在弦外，不能画成两点夹住焦点。\begin{center}\begin{tikzpicture}[font=\small,>=stealth]
\draw[rounded corners=3pt,gray!45] (0,0) rectangle (14.2,15);
\node at (7.1,14.35) {异支、同向：两点在焦点同侧，较近段的倒数更大};\begin{scope}[shift={(4.5,9.0)},scale=0.37]
\begin{scope}\clip (-10,-4) rectangle (10,9);
\draw[teal!75!black,thick,samples=110,domain=-1.9:1.9,variable=\t] plot ({2*cosh(\t)},{4*sinh(\t)});
\draw[teal!75!black,thick,samples=110,domain=-1.9:1.9,variable=\t] plot ({-2*cosh(\t)},{4*sinh(\t)});
\end{scope}
\draw[->,gray] (-10.2,0)--(10.3,0);
\draw[orange!85!black,very thick] (4.47213595499958,0)--(-4.4721359550,8.0000000000);\draw[blue!75!black,very thick] (4.47213595499958,0)--(2.2360679775,2.0000000000);
\fill (2.2360679775,2.0000000000) circle(6pt) node[above right=7pt] {$A$};\fill (-4.4721359550,8.0000000000) circle(6pt) node[left=14pt] {$B$};
\fill (4.47213595499958,0) circle(6pt) node[below right=9pt] {$F$};
\end{scope}

\node[anchor=west,align=left] at (9.5,10.0)
 {图取 $a=2,b=4$\\[13pt]$AF=3,\quad BF=12$\\[13pt]$A$ 在右支，$B$ 在左支\\[13pt]顺序是 $F-A-B$\\[13pt]右支点 $A$ 更近};

\node at (7.1,7.15) {把这条斜线展开看长度：$BF=AF+AB=3+9=12$};
\draw[orange!85!black,very thick] (2.2,6.0)--(8.2,6.0);
\draw[blue!75!black,very thick] (6.7,6.0)--(8.2,6.0);
\fill (2.2,6.0) circle(2pt) node[above=8pt] {$B$};
\fill (6.7,6.0) circle(2pt) node[above=8pt] {$A$};
\fill (8.2,6.0) circle(2pt) node[above=8pt] {$F$};
\draw[<->,orange!85!black] (2.2,5.45)--(8.2,5.45);
\node[orange!85!black] at (5.2,4.9) {整段 $BF=12$};
\node[anchor=west] at (9.3,6.0) {$AF=3,\ AB=9$};
\node at (7.1,3.6) {$\displaystyle \frac13-\frac1{12}=\frac14=\frac{2a}{b^2}$};
\fill[blue!30] (2.8,2.15) rectangle (7.6,2.65);
\fill[orange!40] (6.4,2.15) rectangle (7.6,2.65);
\draw[gray] (2.8,2.15) rectangle (7.6,2.65);
\draw[orange!85!black,thick] (6.4,2.0)--(7.6,2.8) (6.4,2.8)--(7.6,2.0);
\node[blue!75!black] at (4.6,1.45) {剩下 $1/4$};
\node[orange!85!black,anchor=west] at (8.2,2.4) {减掉 $1/BF=1/12$};
\node at (7.1,.6) {整条色条原长是 $1/AF=1/3$；减去橙色部分，留下固定差。};
\end{tikzpicture}
\end{center}
两支的焦半径关系不同：右支 $x_A-c=(r_A-p)/e$，左支 $x_B-c=(-r_B-p)/e$。代入同向比例，得到
\[
 \frac{r_A-p}{er_A}=\frac{-r_B-p}{er_B}
 \quad\Longrightarrow\quad \frac1{r_A}-\frac1{r_B}=\frac2p>0.
\]
因此右支点 $A$ 更近。若没有规定哪点叫 $A$，倒数差外面保留绝对值。
\end{explanationbox}

\clearpage
\begin{explanationbox}
  \textbf{第五步，最容易混淆的事：焦准左右配对，点却不分左右支。}

  同一条双曲线有两套\textbf{完整搭档}：
  \[
    \begin{array}{lll}
      \text{左焦点 }F_1(-c,0)&\longleftrightarrow&\text{左准线 }\ell_1:x=-a^2/c,\\[3pt]
      \text{右焦点 }F_2(c,0)&\longleftrightarrow&\text{右准线 }\ell_2:x=a^2/c.
    \end{array}
  \]
  \textbf{搭档的是焦点与准线，不是“左支点与左焦点”。}
  对左右任意一支上的点 $P$，任选\textbf{一整套}，都满足
  \[
    \frac{|PF_i|}{d(P,\ell_i)}=e=\frac ca>1\qquad(i=1,2).
  \]

  例如选\textbf{右焦点 $F_2$ 和右准线 $\ell_2$}，过该焦点的直线交左支于 $A$、右支于 $B$：
  \[
    \begin{aligned}
      |AF_2|&=e\,d(A,\ell_2)=e\left(\frac{a^2}{c}-x_A\right)
          &&(x_A\le -a),\\[3pt]
      |BF_2|&=e\,d(B,\ell_2)=e\left(x_B-\frac{a^2}{c}\right)
          &&(x_B\ge a).
    \end{aligned}
  \]
  看似两种公式，\textbf{其实只有一条}：
  \[
     \boxed{|PF_2|=e\left|x-\frac{a^2}{c}\right|=|ex-a|}.
  \]
  只是两支相对右准线的位置不同，去绝对值时正负不同。
  \textbf{做 C057 时始终固定右焦点--右准线，就无需走到左支时临时改焦点。}
\end{explanationbox}
